% Erzeugt demonstrationsarbeit-sap-btp.pdf:
% ein vorangestelltes Hinweisblatt und auf jeder Seite der Arbeit
% einen Vermerk, dass es sich um eine KI-Demonstrationsarbeit handelt.
% Eingabe: quelle/thesis-original.pdf (nicht im Repositorium).
\documentclass[11pt,a4paper]{article}
\usepackage[utf8]{inputenc}
\usepackage[T1]{fontenc}
\usepackage[ngerman]{babel}
\usepackage{lmodern}
\usepackage[margin=25mm]{geometry}
\usepackage[autostyle=true,german=quotes]{csquotes}
\usepackage{listings}
\lstset{basicstyle=\ttfamily\small,breaklines=true,frame=single,
  framesep=3mm,columns=fullflexible,
  literate={ö}{{\"o}}1 {ä}{{\"a}}1 {ü}{{\"u}}1 {ß}{{\ss}}1
           {Ö}{{\"O}}1 {Ä}{{\"A}}1 {Ü}{{\"U}}1}
\usepackage{pdfpages}
\usepackage{tikz}
\usepackage[hidelinks]{hyperref}
\pagestyle{empty}
\setlength{\parindent}{0pt}
\setlength{\parskip}{0.6em}

\newcommand{\vermerk}{%
  \begin{tikzpicture}[remember picture,overlay]
    \node[anchor=north,yshift=-5mm,fill=white,fill opacity=0.85,
          text opacity=1,inner sep=3pt,font=\sffamily\footnotesize,
          text=red!60!black]
      at (current page.north)
      {Demonstrationsarbeit -- mit KI erstellt -- keine Prüfungsleistung};
  \end{tikzpicture}}

\begin{document}

\begin{center}
  {\sffamily\bfseries\Large Demonstrationsarbeit\\[0.3em]
  Keine Prüfungsleistung\par}
\end{center}

\vspace{1em}
Die folgende Arbeit mit dem Titel \emph{Generierung von Beispielcode
mit definierten Eigenschaften und Vorhersage der Test-Coverage}
wurde im Oktober 2026 im Rahmen des Beitrags \emph{Sapere, non
scire} mit einem KI-Agenten erstellt (Claude Opus~5.5, Anthropic,
genutzt über Claude Code). Grundlage war allein die unten
wiedergegebene Vorgabe. Sie enthielt den Titel und formale Vorgaben,
aber keine fachliche Gliederung, keine Berechnungen und keinen
Programmcode. Den Prototyp \texttt{ppgen}, die Versuche und ihre
Auswertung hat der Agent ebenfalls selbst entwickelt und ausgeführt;
sie liegen im Ordner \texttt{prototyp} dieses Repositoriums.

Die Arbeit wurde \textbf{nie als Prüfungsleistung eingereicht und
nicht bewertet}. Sie dient als Demonstrationsmaterial für die Frage,
was eine schriftliche Ausarbeitung über die Kompetenz ihres
Abgebenden aussagt, wenn ein KI-Agent sie vollständig erstellen kann.

Die Angaben auf dem folgenden Titelblatt -- Hochschule, Fakultät,
Studiengang, angestrebter Grad und Betreuung -- stammen aus der
verwendeten LaTeX-Vorlage für Abschlussarbeiten. Sie beschreiben
kein reales Prüfungsverfahren; die Betreuenden sind Platzhalter
(\enquote{N.\,N.}). Als Autor ist das Modell eingetragen,
das die Arbeit erstellt hat. Die Vorgabe nennt als
Autorenbezeichnung \enquote{GPT-6.1 Sol (Hoch)}; erstellt wurde die
Arbeit jedoch mit Claude Opus~5.5.

Die Arbeit enthält die \textbf{eidesstattliche Erklärung} der
Hochschulvorlage, unterschrieben mit der Modellbezeichnung. Sie hat
keinerlei rechtliche Bedeutung: Ein KI-System kann keine solche
Erklärung abgeben. Das Vorwort der Arbeit weist selbst darauf hin.

Wer diese Arbeit oder Teile davon als eigene Prüfungsleistung
einreicht, täuscht.

\medskip
\textbf{Vollständige Vorgabe (unverändert):}
\begin{center}
\begin{minipage}{0.9\linewidth}
\lstinputlisting{prompt.txt}
\end{minipage}
\end{center}

\clearpage
\includepdf[pages=-,pagecommand={\thispagestyle{empty}\vermerk}]%
  {quelle/thesis-original.pdf}

\end{document}
